\documentclass[11pt]{article}
\usepackage[margin=1in]{geometry}
\usepackage{amsmath,amssymb,amsthm,hyperref,enumitem}
\hypersetup{colorlinks=true,urlcolor=blue,linkcolor=blue}
\setlength{\emergencystretch}{1em}
\newtheorem{theorem}{Theorem}
\title{Hilbert coefficients of the square of the maximal ideal}
\author{Principia Math In Progress}
\date{22 September 2026}
\begin{document}
\maketitle
\noindent\textbf{MathDB 360642. Full intended formula proved;
zero new campaign credit.}
The coefficient identity is already a direct specialization of classical
Gegenbauer formulas. We record that prior implication explicitly and
supply elementary algebraic and combinatorial proofs. Neither the
alternative proof nor the source's index correction earns new credit.

\section*{Task 1: the complete intended statement}
Let \(S_n=k[x_1,\ldots,x_n]\), where \(k\) is any field and \(n\geq1\),
and put \(I=(x_1,\ldots,x_n)^2\). Use the source's defining convention
\[
 \ell(S_n/I^s)=\sum_{j=0}^n(-1)^j e_j(I)
                    \binom{s+n-j-1}{n-j}\qquad(s\gg0).       \tag{1}
\]
\begin{theorem}
For every \(n\geq1\),
\[
 e_j(I)=
 \begin{cases}
 2^{n-2j}\binom{n-j}{j},&0\leq j\leq\lfloor n/2\rfloor,\\
 0,&\lfloor n/2\rfloor<j\leq n.
 \end{cases}                                                \tag{2}
\]
In particular the claimed nonzero coefficients are strictly positive.
The corresponding identity, valid for every integer \(s\geq1\), is
\[
 \binom{2s+n-1}{n}
 =\sum_{j=0}^{\lfloor n/2\rfloor}(-1)^j2^{n-2j}
        \binom{n-j}{j}\binom{s+n-j-1}{n-j}.                  \tag{3}
\]
\end{theorem}
The current preprint's Conjecture 3 and the published Conjecture 3.1
state the intended coefficient formula (2), but omit the final minus
one in the upper entry of the last binomial in (3). Their own example
\(n=4,s=5\) uses the corrected convention. The uncorrected identity
already fails at \(n=1\), giving \(2s=2s+2\). This typographical issue
is separate from proving (2) for every dimension.

\section*{An elementary proof in every dimension}
The monomials of total degree less than \(2s\) form a basis of
\(S_n/I^s\), so stars and bars gives its length as the left side of (3).
In formal power series, for a fixed positive integer \(s\),
\[
 (1-t)^{-2s}=(1-2t+t^2)^{-s}
      =\sum_{r\geq0}\binom{s+r-1}{r}(2t-t^2)^r.             \tag{4}
\]
All substitutions are valid formally because \(2t-t^2\) has zero
constant term. Its \(r\)-th power contributes to degree \(n\) only
when \(r=n-j\) and \(0\leq j\leq\lfloor n/2\rfloor\); that
coefficient is \((-1)^j2^{n-2j}\binom{n-j}{j}\).
The coefficient of \(t^n\) on the left of (4) is
\(\binom{2s+n-1}{n}\), proving (3).

For fixed \(n\), the polynomials
\(\binom{s+n-j-1}{n-j}\), \(0\leq j\leq n\), have distinct
degrees \(n-j\) and nonzero leading coefficients. They form a basis
of the polynomials of degree at most \(n\). Comparing (1) and (3)
therefore gives every coefficient in (2), including the vanishing
ones. This argument concerns integer dimensions and lengths and is
valid independently of the characteristic of \(k\).

\section*{A combinatorial proof by cancellation}
Consider \(s\) ordered, possibly empty words made from tiles \(A,B,D\).
Tiles \(A,B\) each have size one and positive sign; \(D\) has size
two and negative sign. Give a tuple of words total size \(n\) and
sign \((-1)^{\#D}\). When it has \(j\) tiles \(D\), it has
\(r=n-j\) tiles altogether. Choose the \(j\) domino locations,
color the other tiles \(A\) or \(B\), and distribute the ordered
tile string among the \(s\) words by a weak composition of \(r\).
Its signed count is precisely the \(j\)-th summand of (3).

Scan the words in order and the tiles within each word from left
to right. Find the first occurrence of either a tile \(D\) or an
adjacent pair \(AB\) within one word. Replace
\[
                        D\ \longleftrightarrow\ AB.         \tag{5}
\]
If no such occurrence exists, leave the tuple fixed. The replacement
preserves size and flips the sign. It is an involution: the unchanged
prefix has no earlier occurrence, and expanding \(D\) starts with
\(A\), so no new \(AB\) starts just before the replacement.
Conversely contracting \(AB\) to \(D\) creates no earlier occurrence.
The same place is consequently selected on applying (5) again.
Word boundaries are never crossed.

Every fixed tuple has no domino and no \(AB\) in any word. Thus each
word is uniquely \(B^{b_i}A^{a_i}\), with \(a_i,b_i\geq0\).
The fixed tuples are in bijection with the weak compositions
\[
                  a_1+b_1+\cdots+a_s+b_s=n.
\]
They all have positive sign, and their number is
\(\binom{n+2s-1}{n}\). Cancellation proves (3) for every \(n,s\),
and the polynomial-basis argument above completes (2).

\section*{The classical identity and counting decision}
This result is already implied by the two classical identities
\[
 C_n^{(s)}(x)=\sum_{j=0}^{\lfloor n/2\rfloor}
       \frac{(-1)^j(s)_{n-j}}{j!(n-2j)!}(2x)^{n-2j},\qquad
 C_n^{(s)}(1)=\frac{(2s)_n}{n!}.                            \tag{6}
\]
Here \((a)_r\) is the rising factorial. The first formula is NIST
DLMF 18.5.10; the second follows immediately from 18.5.9 at \(x=1\),
where the hypergeometric argument is zero. Since
\[
 \frac{(s)_{n-j}}{j!(n-2j)!}
       =\binom{n-j}{j}\binom{s+n-j-1}{n-j},
\]
their specialization is exactly (3). DLMF 18.12.4 also gives the
generating function behind (4). The relevant equations and parameter
specializations were inspected directly. The full DLMF chapter and
its cited books were not read.

The complete current and published problem sources still call this
a conjecture. Nevertheless, the existing full identity above already
implies every intended coefficient, so we conservatively exclude
this package from new campaign credit. No claim of a new binomial
identity or historical priority is made for either elementary proof.

\section*{Tasks 2--5: four proposed follow-ups}
\begin{enumerate}[label=\textbf{Task \arabic*:},start=2,leftmargin=*]
\item Extend the explicit coefficient extraction to every ideal
\((x_1,\ldots,x_n)^d\), with \(d\geq3\).
\item Find a cancellation model for those higher-power coefficient
formulas, with fixed objects explaining their signs.
\item Derive exact Hilbert series for powers of ideals generated by
all squares and the edges of an arbitrary finite graph.
\item Determine which such graph ideals have nonnegative Hilbert
coefficients in every dimension.
\end{enumerate}
These are proposed mathematical directions, not additional solutions.

\section*{Sources and verification}
R. Fr\"oberg, \emph{Hilbert coefficients of quadratic algebras},
\href{https://arxiv.org/abs/2311.02860}{arXiv:2311.02860v1},
6 November 2023, Conjecture 3, p.~5; all seven pages read.
Published as \emph{Journal of Experimental Mathematics} 1(2) (2025),
315--321, \href{https://doi.org/10.56994/JXM.001.002.007}{DOI
10.56994/JXM.001.002.007}, Conjecture 3.1, p.~319; all seven article
pages and both front-matter pages read.

NIST Digital Library of Mathematical Functions:
\href{https://dlmf.nist.gov/18.5.E9}{18.5.9},
\href{https://dlmf.nist.gov/18.5.E10}{18.5.10}, and
\href{https://dlmf.nist.gov/18.12.E4}{18.12.4}.
Exact equations and their displayed context checked 22 September 2026.

Supplementary exact checks independently recover the Hilbert
coefficients from lengths, enumerate the cancellation involution on
small word tuples, and test the displayed index correction.
Finite checks supplement the proofs; they do not prove the infinite
claims. Review by the authoring agent only, without independent peer
review or a proof assistant.
\end{document}
